% Verbatim from /Users/alingalatan/Downloads/cubic_chaos_partial_results (2).zip:large3_spectral_reduction.tex (16 November 2026); GPT (a ChatGPT session of the user), relayed by the user.
% Intended filepath: /workspace/scratch/3efc7fe714fa/large3_spectral_reduction.tex
\documentclass{article}
\usepackage{amsmath,amssymb,amsthm}
\begin{document}
\section*{An explicit coordinate reduction, including the mixed terms}
All results in this note are rigorous. The final limitation is explicit:
the reduction is not a proof of the general sharp $L^1$ bound.

Let $F=I_3(f)$ be a normalized real third-chaos element,
$6\|f\|^2=1$. Let $A=f\otimes_2f$ be its positive semidefinite
one-coordinate contraction matrix. Thus $\operatorname{tr}A=1/6$.
Given $\theta>0$, let $V$ be the span of eigenvectors with eigenvalues
strictly greater than $\theta$, and let $P$ and $Q=1-P$ be its orthogonal
projectors. Then
\[
 r:=\dim V\le\frac1{6\theta}.
 \tag{1}
\]
Let $f_0=Q^{\otimes3}f$, $R=I_3(f_0)$, and $v=\mathbb ER^2$.
For $u\in V^\perp$,
\[
 \langle (f_0\otimes_2 f_0)u,u\rangle
 =\|Q^{\otimes2}f(u,\cdot,\cdot)\|^2
 \le\|f(u,\cdot,\cdot)\|^2=\langle Au,u\rangle
 \le\theta\|u\|^2.
\]
Writing $A_0=f_0\otimes_2f_0$, we therefore have $0\le A_0\le\theta I$,
$\operatorname{tr}A_0=v/6$, and the exact contraction formula gives
\[
 \kappa_4(R)\le3240\operatorname{tr}(A_0^2)
 \le540\theta v.
 \tag{2}
\]
If $v>0$, division by $v^2$ yields
\[
 \kappa_4(R/\sqrt v)\le\frac{540\theta}{v}.
 \tag{3}
\]
Consequently, for arbitrary $\epsilon,\eta>0$, choose
$\theta=\epsilon\eta^2/540$. Then
\[
 \boxed{r\le\frac{90}{\epsilon\eta^2},\qquad
 \|R\|_2<\eta\quad\hbox{or}\quad
 \kappa_4(R/\|R\|_2)\le\epsilon.}
 \tag{4}
\]
Here $\eta$ is an $L^2$-norm threshold. If it is instead a variance
threshold, replace $\eta^2$ throughout by $\eta$.

\subsection*{What remains after removing the pure-tail cubic}
Write the independent Gaussian vectors on $V$ and $V^\perp$ as $X$ and $Y$.
The exact orthogonal decomposition is
\[
 F=P_3(X)+\sum_{|\alpha|=2}\mathrm{He}_\alpha(X)L_\alpha(Y)
      +\sum_{i=1}^r X_iQ_i(Y)+R(Y),
 \tag{5}
\]
where $P_3\in H_3(X)$, $L_\alpha\in H_1(Y)$, and $Q_i\in H_2(Y)$.
The four displayed pieces have distinct head degrees and are orthogonal.
The tail variables appearing in the three last pieces are generally
dependent across those pieces; no independence assertion is made.

There is a quantitative restriction on the mixed quadratic coefficients.
For a unit basis vector $e_i\in V$, set
$M_i=Qf(e_i,\cdot,\cdot)Q$. Then $Q_i(Y)=3I_2(M_i)$, and for tail $u$,
\[
 \|M_i u\|^2\le\sum_{j,k}f(u,e_j,e_k)^2
 =\langle Au,u\rangle\le\theta\|u\|^2.
\]
Thus $\|M_i\|_{\mathrm{op}}\le\sqrt\theta$, and the quadratic matrix
in $Q_i$ has operator norm at most $3\sqrt\theta$. This alone is not
a joint Gaussian approximation statement, especially when $r$ grows.

For completeness, large cumulant forces a definite amount of variance in
the structured pieces of (5). Put $z=\operatorname{tr}(PA)$ and
$\delta=\kappa_4(F)$. Since
\[
 \frac\delta{3240}\le\operatorname{tr}(A^2)
 \le\frac z6+\frac\theta6,
\]
we have $z\ge\delta/540-\theta$. Also $1-v\ge6z$, because tensor
entries whose first index lies in $V$ are among the entries having at least
one index in $V$. Therefore
\[
 1-v\ge\frac\delta{90}-6\theta.
 \tag{6}
\]
For $\theta\le\delta/1080$ this gives $1-v\ge\delta/180$.

\subsection*{A concrete mixed obstruction to a naive reduction}
Let $X,Y_1,\ldots,Y_N$ be independent standard Gaussians and set
\[
 Q_N(Y)=\frac1{\sqrt{2N}}\sum_{j=1}^N(Y_j^2-1),\qquad
 F_N=XQ_N(Y).
\]
Then $F_N\in H_3$, $\mathbb EF_N^2=1$, and
\[
 \kappa_4(F_N)=3\bigl(3+12/N\bigr)-3=6+36/N\longrightarrow6.
\]
The eigenvalue of $A_N$ in the $X$ direction is $1/18$, while every $Y_j$
eigenvalue is $1/(9N)$. Hence a fixed cutoff $0<\theta<1/18$ selects
only $X$ for sufficiently large $N$. In (5), $P_3=0$, all $L_\alpha=0$,
and $R=0$, while the entire $F_N$ remains in the mixed term $XQ_N(Y)$.
Thus controlling the pure-tail remainder does not remove the mixed terms.
This family itself satisfies the desired inequality, since
\[
 \mathbb E|F_N|=\sqrt{2/\pi}\,\mathbb E|Q_N|
 \le\sqrt{2/\pi}.
\]
The unresolved step is a sharp estimate for general sums of all the
pieces in (5), retaining their dependence. Treating $R$ or the $Q_i$ as
independent Gaussian variables without a joint approximation theorem would
be heuristic and is not part of this reduction.
\end{document}
